\documentclass{article}

\usepackage{../handouts}

\title{CSCI 432 Handout 11: Greedy Algorithms}
\author{Name: \rule{3in}{0.15mm}}

\begin{document}
\maketitle

\section*{Class Scheduling}

Consider the problem of class scheduling. There are $n$ classes available, and
you want to take the maximum number of classes.  Each class has a start time and
a finish time:
$$ C = \{ c_i =(s_i,f_i) \}_{i=1}^n.$$
What strategies allow you to take the most classes? For this, we will consider
several different strategies. We will call a strategy optimal if it is one that
maximizes the number of classes that you can take.

For the strategies listed
below, are they optimal or not? If not, provide a counter-example to prove that
it is not optimal.

\begin{enumerate}
    \item Starts first: Choose the class that starts first, eliminate any
        conflicts with that. Among the remaining, choose the class that starts
        first and repeat this process until no classes are remaining.

        \textbf{Answer}
        \vspace{1in}

    \item Starts last.

        \textbf{Answer}
        \vspace{1in}

    \item Finishes first.

        \textbf{Answer}
        \vspace{1in}

    \item Finishes last.

        \textbf{Answer}
        \vspace{1in}

    \item Longest duration.

        \textbf{Answer}
        \vspace{1in}

    \item Shortest duration.

        \textbf{Answer}
        \vspace{1in}

    \item Bonus: Randomized selection (here, find an example that shows where
        random selection will have a low probability of finding an optimal
        solution).

        \textbf{Answer}
        \vspace{1in}

\end{enumerate}

\newpage
\section*{Hoffman Codes}

Consider the following code (generated using CatChat\footnote{Prompt was ``latex
table with a prefix-free code for the alphabet.'' The answer had a formatting
problem, but was easily fixable. However, the code given was not prefix free, so
I had to fix that.}:

\begin{table}[ht]
    \centering
    \caption{Prefix‑free binary code for the English alphabet}
    \label{tab:prefixfree}
    \begin{tabular}{cl}
        \toprule
        \textbf{Letter} & \textbf{Codeword} \\
        \midrule
        A & 00000 \\
        B & 11111000 \\
        C & 111100 \\
        D & 1111101 \\
        E & 01100 \\
        F & 1111010 \\
        G & 1111011 \\
        H & 1100 \\
        I & 0010 \\
        J & 111111100 \\
        K & 111111101 \\
        L & 111001 \\
        M & 11111001 \\
        N & 0100 \\
        O & 00010 \\
        P & 0011 \\
        Q & 0111 \\
        R & 1101 \\
        S & 111010 \\
        T & 0101 \\
        U & 111000 \\
        V & 00001 \\
        W & 111011 \\
        X & 01101 \\
        Y & 00011 \\
        Z & 11111111 \\
        \bottomrule
    \end{tabular}
\end{table}

\begin{enumerate}\setcounter{enumi}{2}

    \item Can you decode the following message:

            0110011111001000000010111001%
            1111100101100%
            1111010000101101%
            11111000000100100111000111010

        \textbf{Answer}
        \vspace{1in}

    \pagebreak
    \item For this code, draw a binary tree, withe the encoded characters stored at the leaf
        nodes.

        \textbf{Answer}
        \vspace{1in}

    \item For the Huffman strategy of ``Merge the two least frequent letters and
        recurse", draw a the Huffman tree for the following frequency count:

        \begin{flushleft}
        \begin{table}[ht]
            \caption{Random frequency count for letters in an imaginary text
            that loves weird letters.}
            \label{tab:prefixfree}
            \begin{tabular}{cl}
                \toprule
                \textbf{Letter} & \textbf{Codeword} \\
                \midrule
                A &  7\\
                B & 13\\
                C &  2\\
                D & 15\\
                E &  4\\
                F & 12\\
                G &  0\\
                H &  9\\
                I &  5\\
                J & 11\\
                K &  1\\
                L & 14\\
                M &  3\\
                N & 10\\
                O &  6\\
                P &  8\\
                Q & 15\\
                R &  2\\
                S & 13\\
                T &  4\\
                U &  9\\
                V &  0\\
                W & 12\\
                X &  5\\
                Y &  7\\
                Z & 11\\
                \bottomrule
            \end{tabular}
        \end{table}
        \end{flushleft}

\end{enumerate}

\end{document}
